\documentclass{article}

% packages
\usepackage{amsmath, amsthm, thmtools, amsfonts, amssymb, luacode, catchfile, tikzducks, hyperref, ifthen}
\ifcsname c@kobocompile\endcsname
	\usepackage[a5paper, total={1072pt, 1448pt}, margin=10pt, includeheadfoot]{geometry} % set page margins
\else
	\usepackage[a4paper, margin=50pt, includeheadfoot]{geometry}
\fi
\usepackage[shortlabels]{enumitem}
\usepackage[skip=3pt, indent=0pt]{parskip}

% language
\usepackage[bidi=basic, layout=tabular, provide=*]{babel}
\ifcsname c@english\endcsname
	\babelprovide[main, import]{english}
\else
	\babelprovide[main, import]{hebrew}
	\babelprovide{rl}
\fi
%\babelfont{rm}{Libertinus Serif}
\babelfont{rm}[Renderer=Harfbuzz]{Libertinus Serif}
\babelfont{sf}{Libertinus Sans}
\babelfont{tt}{Libertinus Mono}

% style
\AddToHook{cmd/section/before}{\clearpage}	% Add line break before section
\linespread{1.3}
\setcounter{secnumdepth}{0}		% Remove default number tags from sections, this won't do well with theorems
\AtBeginDocument{\setlength{\belowdisplayskip}{3pt}}
\AtBeginDocument{\setlength{\abovedisplayskip}{3pt}}
\graphicspath{ {../images/} }

% operators
\DeclareMathOperator\cis{cis}
\DeclareMathOperator\Sp{Sp}
\DeclareMathOperator\tr{tr}
\DeclareMathOperator\im{Im}
\DeclareMathOperator\re{Re}
\DeclareMathOperator\diag{diag}
\DeclareMathOperator*\lowlim{\underline{lim}}
\DeclareMathOperator*\uplim{\overline{lim}}
\DeclareMathOperator\rng{rng}
\DeclareMathOperator\Sym{Sym}
\DeclareMathOperator\Arg{Arg}
\DeclareMathOperator\Log{Log}
\DeclareMathOperator\dom{dom}
\DeclareMathOperator\supp{Supp}
\DeclareMathOperator\var{Var}
\DeclareMathOperator\cov{Cov}

% commands
%\renewcommand\qedsymbol{\textbf{מש''ל}}
%\renewcommand\qedsymbol{\fbox{\emoji{lizard}}}
\newcommand{\Aa}[0]{\mathcal{A}}
\newcommand{\Bb}[0]{\mathcal{B}}
\newcommand{\CC}[0]{\mathbb{C}}
\newcommand{\Cc}[0]{\mathcal{C}}
\newcommand{\EE}[0]{\mathbb{E}}
\newcommand{\FF}[0]{\mathbb{F}}
\newcommand{\Ff}[0]{\mathcal{F}}
\newcommand{\Ii}[0]{\mathcal{I}}
\newcommand{\Gg}[0]{\mathcal{G}}
\newcommand{\Ll}[0]{\mathcal{L}}
\newcommand{\Mm}[0]{\mathcal{M}}
\newcommand{\NN}[0]{\mathbb{N}}
\newcommand{\Nn}[0]{\mathcal{N}}
\newcommand{\PP}[0]{\mathbb{P}}
\newcommand{\Pp}[0]{\mathcal{P}}
\newcommand{\QQ}[0]{\mathbb{Q}}
\newcommand{\RR}[0]{\mathbb{R}}
\newcommand{\Rr}[0]{\mathcal{R}}
\newcommand{\Ss}[0]{\mathcal{S}}
\newcommand{\TT}[0]{\mathbb{T}}
\newcommand{\Uu}[0]{\mathcal{U}}
\newcommand{\Vv}[0]{\mathcal{V}}
\newcommand{\Ww}[0]{\mathcal{W}}
\newcommand{\ZZ}[0]{\mathbb{Z}}
\newcommand{\acts}[0]{\circlearrowright}
\newcommand{\explain}[2] {
	\begin{flalign*}
		 && \text{#2} && \text{#1}
	\end{flalign*}
}
\newcommand{\maketitleprint}[0]{ \begin{center}
	%\begin{tikzpicture}[scale=3]
	%	\duck[graduate=gray!20!black, tassel=red!70!black]
	%\end{tikzpicture}	
	\includegraphics[width=6cm]{cover}
\end{center}
}

% theorem commands
\newtheoremstyle{c_remark}
	{}	% Space above
	{}	% Space below
	{}% Body font
	{}	% Indent amount
	{\bfseries}	% Theorem head font
	{}	% Punctuation after theorem head
	{.5em}	% Space after theorem head
	{\thmname{#1}\thmnumber{ #2}\thmnote{ \normalfont{\text{(#3)}}}}	% head content
\newtheoremstyle{c_definition}
	{3pt}	% Space above
	{3pt}	% Space below
	{}% Body font
	{}	% Indent amount
	{\bfseries}	% Theorem head font
	{}	% Punctuation after theorem head
	{.5em}	% Space after theorem head
	{\thmname{#1}\thmnumber{ #2}\thmnote{ \normalfont{\text{(#3)}}}}	% head content
\newtheoremstyle{c_plain}
	{3pt}	% Space above
	{3pt}	% Space below
	{\itshape}% Body font
	{}	% Indent amount
	{\bfseries}	% Theorem head font
	{}	% Punctuation after theorem head
	{.5em}	% Space after theorem head
	{\thmname{#1}\thmnumber{ #2}\thmnote{ \text{(#3)}}}	% head content

\ifcsname c@english\endcsname
	\theoremstyle{plain}
	\newtheorem{theorem}{Theorem}[section]
	\newtheorem{lemma}[theorem]{Lemma}
	\newtheorem{proposition}[theorem]{Proposition}
	\newtheorem*{proposition*}{Proposition}
	%\newtheorem{corollary}[theorem]{אין חלופה עברית}

	\theoremstyle{definition}
	\newtheorem{definition}[theorem]{Definition}
	\newtheorem*{definition*}{Definition}
	\newtheorem{example}{Example}[section]
	\newtheorem{exercise}{Exercise}[section]

	\theoremstyle{remark}
	\newtheorem*{remark}{Remark}
	\newtheorem*{solution}{Solution}
	\newtheorem{conclusion}[theorem]{Conclusion}
	\newtheorem{notation}[theorem]{Notation}
\else
	\theoremstyle{c_plain}
	\newtheorem{theorem}{משפט}[section]
	\newtheorem{lemma}[theorem]{למה}
	\newtheorem{proposition}[theorem]{טענה}
	\newtheorem*{proposition*}{טענה}
	%\newtheorem{corollary}[theorem]{אין חלופה עברית}

	\theoremstyle{c_definition}
	\newtheorem{definition}[theorem]{הגדרה}
	\newtheorem*{definition*}{הגדרה}
	\newtheorem{example}{דוגמה}[section]
	\newtheorem{exercise}{תרגיל}[section]

	\theoremstyle{c_remark}
	\newtheorem*{remark}{הערה}
	\newtheorem*{solution}{פתרון}
	\newtheorem{conclusion}[theorem]{מסקנה}
	\newtheorem{notation}[theorem]{סימון}
\fi

% Questions related commands
\newcounter{question}
\setcounter{question}{1}
\newcounter{sub_question}
\setcounter{sub_question}{1}

\ifcsname c@english\endcsname
	\newcommand{\question}[1][0]{
		\ifthenelse{#1 = 0}{}{\setcounter{question}{#1}}
		\section{Question \arabic{question}}
		\addtocounter{question}{1}
		\setcounter{sub_question}{1}
	}

	\newcommand{\subquestion}[1][0]{
		\ifthenelse{#1 = 0}{}{\setcounter{sub_question}{#1}}
		\subsection{Part \alph{sub_question}}
		\addtocounter{sub_question}{1}
	}
\else
	\newcommand{\question}[1][0]{
		\ifthenelse{#1 = 0}{}{\setcounter{question}{#1}}
		\section{שאלה \arabic{question}}
		\addtocounter{question}{1}
		\setcounter{sub_question}{1}
	}

	\newcommand{\subquestion}[1][0]{
		\ifthenelse{#1 = 0}{}{\setcounter{sub_question}{#1}}
		\subsection{סעיף \localecounter{letters.gershayim}{sub_question}}
		\addtocounter{sub_question}{1}
	}
\fi

% import lua and start of document
\directlua{common = require ('../common')}

\GetEnv{AUTHOR}

% headers
\author{\AUTHOR}
\date\today

\title{פתרון מטלה 1 --- גאומטריה דיפרנציאלית אלמנטרית, 80560}

\begin{document}
\maketitle
\maketitleprint[purple]

\section{שאלה 0}
יהי $(\Pp, \Ll)$ מישור אפיני בגישה סינתטית.

\subquestion{}
נוכיח כי במישור האפיני יש לפחות שלושה ישרים שונים.
\begin{proof}
	נסמן ב־$P, Q, R, S$ את 4 הנקודות הלא קולינאריות שנתון ושמצאנו שקיימות.
	מצאנו במהלך הוכחה ש־$l = \langle P, Q \rangle$ ו־$m = \langle P, R \rangle$ יחד עם $l' \parallel l \ni R$ וכן $m \parallel m' \ni S$ הם ישרים שונים.
\end{proof}

\subquestion{}
נוכיח שלא קיים ישר ללא נקודות.
\begin{proof}
	נניח בשלילה שקיים ישר כזה $l \in \Ll$.
	נניח גם ש־$P, Q, R \in \Pp$ נקודות לא קולינאריות, וכן נסמן $m = \langle P, Q \rangle, n = \langle P, R \rangle$, אז ידוע ש־$m \not\parallel n$.
	נגדיר את המשיקים $l \parallel l' \ni Q$ וכן את $l \parallel l'' \ni R$, אז בהכרח $l' \parallel l''$ וכן $Q \in l' \cap m$ וגם $R \in l'' \cap R$.
	אילו $l \parallel m$ וגם $l \parallel n$ אז נקבל ש־$l' = l''$ ולכן $P, Q, R$ קולינאריות בסתירה, ולכן $l \lnot\parallel m$ וקיימת נקודת חיתוך.
\end{proof}

\subquestion{}
נוכיח כי לכל ישר לפחות שתי נקודות שונות.
\begin{proof}
	יהי ישר $l \in \Ll$, ותהי $P \in l$ נקודה כלשהי שידוע שקיימת מהסעיף הקודם.
	נניח שוב ש־$P, Q, R \in \Pp$ נקודות לא קולינאריות ונגדיר את $m = \langle P, Q \rangle$. \\
	נסמן גם את $n = \langle Q, R \rangle$, אם $n \lnot\parallel l$ אז סיימנו, שכן $Q \in l$ או $R \in l$, לכן נניח ש־$n \parallel l$.
	נגדיר את המשיק $R \in m' \parallel m$, אז $l \not\parallel m'$ בהכרח ולכן יש להם נקודת חיתוך.
\end{proof}

\subquestion{}
נראה כי לכל שני ישרים כמות זהה של נקודות.
\begin{proof}
	נניח ש־$l_0, l_1 \in \Ll$ ישרים מקבילים, ונניח ש־$P_0 \in l_0, P_1 \in l_1$. \\
	אם $l_0 = l_1$ אז סיימנו, לכן נניח ש־$l_0 \ne l_1$, כלומר הנקודות המרכיבות אותם שונות. \\
	נגדיר את הישר $m = \langle P_0, P_1 \rangle$, בהכרח $m \not\parallel l_0, l_1$.
	תהי נקודה $Q \in l_0$, אז נגדיר את $Q \in n \parallel m$, מתקיים $n \not\parallel l_1$ ולכן יש להם נקודת חיתוך $Q' \in l_1$. 

	נניח עתה ש־$Q_0, Q_1 \in l_0$ וכן ש־$Q_0', Q_1' \in l_1$ מתאימות אליה, נראה ש־$Q_0 \ne Q_1 \implies Q_0' \ne Q_1'$.
	נניח ש־$Q_0' = Q_1'$, וכן נניח ש־$Q_0 \ne Q_1$, אז $o = \langle Q_0, Q_0' \rangle, p = \langle Q_1, Q_0' \rangle$ הם ישרים לא מקבילים.
	אבל $o \parallel m \parallel p \implies o \parallel p$ בסתירה.
\end{proof}

\subquestion{}
נראה שלכל שני ישרים נחתכים יש אותה כמות של נקודות.
\begin{proof}
	נפעל באופן זהה לסעיף הקודם.
	אם $l_0, l_1 \in \Ll$ כך ש־$l_0 \cap l_1 = \{ O \}$ אז קיימות נקודות שונות $P_0, P_1 \in \Pp$ כך ש־$P_0 \in l_0, P_1 \in l_1$, מכאן ההוכחה זהה תוך שימוש ב־$m = \langle P_0, P_1 \rangle$.
\end{proof}

\subquestion{}
נסיק שלכל הישרים במישור האפיני יש אותה כמות של נקודות.
\begin{proof}
	נסמן שתי נקודות כלשהן $P, Q \in \Pp$ ואת $l = \langle P, Q \rangle$.
	יהי $m \in \Ll$ ישר כלשהו.
	אם $l \parallel m$ אז יש להם אותו מספר נקודות בהתאם לסעיף ד'.
	אם $l \not\parallel m$ אז מסעיף ה' יש לישרים אותו מספר נקודות.
\end{proof}

\question{}
יהי $\FF$ שדה ויהי $(E, V, t)$ מרחב אפיני. \\
תהי $P \in E$ ו־$(E, +_P, \cdot_P)$ מרחב וקטורי מעל $\FF$ שראשיתו $P$, נסמן כ־$E_P$. \\
נוכיח שההעתקה $v_P : E_P \to V$ היא איזומורפיזם של מרחבים וקטוריים.
\begin{proof}
	נוכיח תחילה ש־$v_P$ העתקה לינארית.
	נניח ש־$\alpha, \beta \in \FF$ וכן ש־$Q, R \in E_P$, אז מתקיים,
	\begin{align*}
		v_P(\alpha Q +_P \beta R)
		& = v(P, \alpha Q +_P \beta R) \\
		& = \alpha Q +_P \beta R - P \\
		& = \alpha \cdot_P + Q \beta \cdot_P R - 2P \\
		& = \alpha (Q - P) + P + \beta (R - P) + P - 2P \\
		& = \alpha v_P(Q) + \beta v_P(R)
	\end{align*}
	תוך שימוש בהגדרות המופיעות בסיכום.

	ראינו ש־$v_P$ היא העתקה הפיכה בהרצאה ולכן היא איזומורפיזם של מרחבים וקטוריים.
\end{proof}

\question{}
\subquestion{}
יהיו $W, W' \le V$ תתי־מרחבים אפיניים ו־$P, Q \in E$ נקודות.
נראה ישירות כי מתקיים,
\[
	P + W
	= Q + W'
	\iff W = W' \land Q - P \in W
\]
\begin{proof}
	נראה תחילה ש־$P + W = Q + W \iff Q - P \in W$. \\
	נניח ש־$P + W = Q + W$, אז,
	\[
		P + W - Q
		= Q + W - Q
		= \{ w + Q - Q \mid w \in W \}
		= \{ w \mid w \in W \}
		= W
	\]
	ולכן בפרט $P + Q \in W$. \\
	נניח ש־$Q - P \in W$, אז מתקיים,
	\[
		P + W
		= P - Q + Q + W
		= Q + W
	\]
	ונסיק את הטענה הראשונה.

	עתה נראה שאם $R + W = R + W'$ אז $W = W'$.
	יהי $w \in R + W$, אז קיים $w' \in W$ כך ש־$w = R + w'$, וידוע כי גם $w \in W'$ ולכן קיים גם $w'' \in W'$ כך ש־$R + w'' = w$.
	אז בהתאם מתקיים $w'' = w'$ ולכן $W = W'$.

	עתה נעבור להוכחת טענת השאלה, נניח ש$P + W = Q + W'$.
	מתקיים,
	\[
		P + W = Q + W'
		\iff P - Q + W = Q - Q + W'
		\iff P - P + W = Q - P + W'
	\]
	ולכן נובע ש־$W = W'$ ומטענת העזר הראשונה נובע שגם $P - Q \in W$.

	נעבור לכיוון השני ונניח ש־$W = W'$ וכן ש־$P - Q \in W$.
	אז הטענה נובעת ישירות מטענת העזר הראשונה.
\end{proof}

\subquestion{}
תהי $F = P + W \le E$ תת־יריעה עבור $P \in E$ ו־$W \le V$. \\
נראה שלכל $Q \in F$ מתקיים $F = Q + W$.
\begin{proof}
	נבחין כי $Q \in F \iff Q = P + w$ עבור $w \in W$.
	בהתאם $Q - P = w \in W$ ולכן מסעיף הקודם נובע $Q + W = P + W = F$.
\end{proof}

\question{}
\begin{definition}
	יהיו $(F_1, W_1)$ ו־$(F_2, W_2)$ שתי תת־יריעות. \\
	נאמר שהן מקבילות אם $W_1 = W_2$ ונסמן $F_1 \parallel F_2$.
\end{definition}
נוכיח את משפט אוקלידס: לכל $F \subseteq E$ תת־יריעה ו־$P \in E \setminus F$ קיימת $P \in F' \subseteq E$ תת־יריעה יחידה מקבילה ל־$F$.
\begin{proof}
	נניח שמתקיים $F = Q + W$ עבור $Q \in E$ ו־$W \le V$. \\
	בהתאם $F' = P + W$ היא תת־יריעה של $E$, מהגדרתה היא מקבילה ל־$F$.

	נותר אם כן להוכיח יחידות, נניח שגם $F'' \le E$ תת־יריעה מקבילה ל־$F$, ונסיק ש־$F'' = P + W = F'$ מהשאלה הקודמת.
\end{proof}

\question{}
תהיינה תת־יריעות $F_1, F_2 \le E$. \\
נראה כי $F_1 \parallel F_2 \iff \exists v \in V,\ F_1 + v = F_2$ וכן שאם $F \le E$ אז לכל $v, w \in V$ מתקיים $F + v = F + w \iff v - w \in W$.
\begin{proof}
	נסמן $F_1 = P_1 + W_1, F_2 = P_2 + W_2$. \\
	נניח ש־$F_1 + v = F_2$ עבור $v \in V$, אז מתקיים $P_1 + W_1 + v = P_2 + W_2$ ולכן נובע ש־$W_1 = W_2$ (משאלה 3) ובהתאם $F_1 \parallel F_2$.

	נניח ש־$F_1 \parallel F_2$, כלומר $W_1 = W_2$.
	נסמן $v = P_1 - P_2 \in V$, אז מתקיים $P_1 + W_1 + v = P_1 + P_2 - P_1 + W_1 = P_2 + W_2$, כלומר $F_1 + v = F_2$.

	נניח עתה ש־$F \le E$ ויהיו $v, w \in V$.
	נניח גם ש־$F = P + W$. \\
	אם $F + v = F + w$ אז $P + v + W = P + w + W$ ולכן $(P - v) - (P - w) \in W$. \\
	נניח בכיוון ההפוך ש־$v - w \in W$, אז תת־היריעות מקבילות ובהכרח $F + v = F + w$.
\end{proof}

\question{}
תהי $S \subseteq E$ ונגדיר,
\[
	\langle S \rangle
	= \bigcap_{S \subseteq F \le E} F
\]
כלומר נאמר ש־$S$ יוצרת את תת־היריעה $\langle S \rangle$. \\
נוכיח כי $\langle S \rangle \le E$ וכן שהיא מינימלית ביחס ההכלה ביחס לתת־היריעות המכילות את $S$.
\begin{proof}
	אם $S = \emptyset$ אז היא תת־יריעה והיא מינימלית ביחס ההכלה, לכן נניח ש־$S \ne \emptyset$, ותהי $P \in S$ נקודה כלשהי.
	מתקיים $P \in F$ לכל $S \subseteq F \le E$ ולכן $F = P + W$ עבור $W \le V$.
	נוכל אם כן להסיק שמתקיים,
	\[
		\langle S \rangle
		= P + \bigcap_{S - P \subseteq W \le V} W
	\]
	כלומר $\langle S \rangle = P + \langle S - P \rangle$, כאשר $\langle S - P \rangle$ הוא תת־מרחב וקטורי הנוצר על־ידי $S - P$.
	נסמן $W' = \langle S - P \rangle$ ונסיק ש־$S = P + W' \le E$. \\
	נשאר להוכיח שתת־יריעה זו היא מינימלית ביחס ההכלה, אך טענה זו נובעת ישירות ממינימליות ביחס להכלה של תת־מרחב וקטורי נוצר.
\end{proof}

\end{document}
